% Paste into your preamble after loading amsmath and feynman-fun.
% Process: gamma(k) + phi(p) -> gamma(k') + phi(p').
% Slots: upper left, lower left, upper right, lower right.
\pgfkeys{/corrdiag/sqed-compton/.style={
  scale=0.70,
  field={},
  blob=false,
  external-legs=true,
  external-vertices=false,
  center-vertex=true,
  center-vertex-style=dot,
  fourpoint-xspan=1.60,
  fourpoint-yspan=1.10,
  leg-styles={pho,chasca,pho,antsca},
  propagator-arrow-size=4pt,
  leg-labels={\epsilon_\mu^{\mathrm{in}}(k),\phi(p),
              {\epsilon_\nu^{\mathrm{out}}}^{*}(k'),\phi(p')},
  show-momenta=true,
  momentum-label-size=\scriptsize,
  momentum-label-fill=none,
  momentum-label-inner-sep=0.5pt,
  momentum-arrow-size=2pt,
  external-momentum-start=0.22,
  external-momentum-end=0.58,
  external-momentum-label-fraction=0.38,
  momentum-label-gap=0,
  % Empty labels retain continuous external momentum arrows.
  momentum-labels={,,,},
  external-momentum-flow=left-in-right-out
}}

% s-channel: scalar momentum p+k flows from left to right.
\newcommand{\ComptonS}[1][]{%
  \STreeChannelCorr[
    sqed-compton,
    channel-vertex-xspan=0.65,
    channel-internal-line=chasca,
    channel-momentum=p+k,
    channel-momentum-angle=90,
    channel-momentum-distance=9pt,
    channel-momentum-label-position=above,
    channel-momentum-label-yshift=2pt,
    #1
  ]%
}

% u-channel: scalar momentum p-k' flows from bottom to top.
% Reversing start/end reverses the separate momentum arrow.
\newcommand{\ComptonU}[1][]{%
  \UTreeChannelCorr[
    sqed-compton,
    channel-vertex-yspan=0.72,
    channel-internal-line=antsca,
    channel-momentum=p-k',
    channel-momentum-start=0.76,
    channel-momentum-end=0.24,
    channel-momentum-angle=-90,
    channel-momentum-distance=10pt,
    channel-momentum-label-position=left,
    #1
  ]%
}

% Local two-photon/two-scalar (seagull) vertex.
\newcommand{\ComptonContact}[1][]{%
  \FourPointCorr[
    sqed-compton,
    topology=contact,
    contact-pairing=central,
    #1
  ]%
}
